%! Inverses of Exponential Functions — Unit 5, Lesson 5.2 — Group Activity
\documentclass[10pt]{article}
\usepackage{precalculus-article}
\usepackage{precalculus-boxes}
\newcommand{\TallMath}[1]{$\displaystyle #1\rule[-1.4em]{0pt}{3.2em}$}

\begin{document}

\pageheader{Unit 5, Lesson 5.2}{Group Activity}
\namepartnerperiod

\vspace{0.1in}

\begin{scenariobox}[A Biologist's Inverse Problem]{sky}
A scientist models bacterial growth using $P(t) = 5^t$, where $P$ is the population count
and $t$ is time in hours. She needs the \emph{inverse} of this function to determine the time
elapsed given a current population count. Your group will construct and verify this inverse.
\end{scenariobox}

\vspace{0.1in}

\begin{tcolorbox}[colback=white, colframe=black!40,
    title=\textbf{Tier R --- Remediate},
    fonttitle=\bfseries, arc=2mm, left=3mm, right=3mm, top=2mm, bottom=2mm]

\textbf{Work with $g(x) = 4^x$.}

\begin{enumerate}[leftmargin=1.8em, itemsep=10pt]

\item Complete the table for $g(x) = 4^x$:

\vspace{4pt}
\begin{center}
\renewcommand{\arraystretch}{1.8}
\begin{tabular}{|c|c|}
\hline
\rowcolor{sky}
$x$ & $g(x) = 4^x$ \\
\hline
$0$ & \blank{1.8cm} \\
\hline
$1$ & \blank{1.8cm} \\
\hline
$2$ & \blank{1.8cm} \\
\hline
$3$ & \blank{1.8cm} \\
\hline
\end{tabular}
\end{center}

\item Swap the columns to make a table for the \emph{inverse} function $f(x) = \log_4 x$:

\vspace{4pt}
\begin{center}
\renewcommand{\arraystretch}{1.8}
\begin{tabular}{|c|c|}
\hline
\rowcolor{sky}
$x$ & $f(x) = \log_4 x$ \\
\hline
\blank{1.8cm} & $0$ \\
\hline
\blank{1.8cm} & $1$ \\
\hline
\blank{1.8cm} & $2$ \\
\hline
\blank{1.8cm} & $3$ \\
\hline
\end{tabular}
\end{center}

\item Using your table, evaluate: \quad $f(16) = \log_4 16 = $ \blank{2cm} \qquad $f(1) = \log_4 1 = $ \blank{2cm}

\vspace{6pt}

\item Verify the inverse: compute $g(f(16)) = 4^{\,f(16)} = $ \blank{2cm}. Does this equal 16? \blank{1.5cm}

\end{enumerate}
\end{tcolorbox}

\vspace{0.1in}

\begin{tcolorbox}[colback=white, colframe=black!40,
    title=\textbf{Tier A --- Approaching Proficiency},
    fonttitle=\bfseries, arc=2mm, left=3mm, right=3mm, top=2mm, bottom=2mm]

\textbf{Return to the scientist's model: $P(t) = 5^t$.}

\begin{enumerate}[leftmargin=1.8em, itemsep=10pt]

\item Write the inverse function $t(P)$ that gives hours elapsed given a population $P$.

\vspace{4pt}
$t(P) = $ \blank{4cm}

\item Verify using composition. Show both directions:

\vspace{4pt}
$P(t(P)) = 5^{\,t(P)} = $ \blank{4cm} \quad $t(P(t)) = \log_5(5^t) = $ \blank{4cm}

\item How many hours until the population reaches 125? Express your answer in logarithmic notation, then evaluate.

\vspace{4pt}
$t(125) = $ \blank{3cm} $= $ \blank{2cm} hours

\item How many hours until the population reaches 625? Express in logarithmic notation.

\vspace{4pt}
$t(625) = $ \blank{3cm} $= $ \blank{2cm} hours

\item Explain in words: how does the inverse function help the scientist solve a problem
      that the original model $P(t) = 5^t$ cannot answer directly?

\writelines{2}

\end{enumerate}
\end{tcolorbox}

\vspace{0.1in}

\begin{tcolorbox}[colback=white, colframe=black!40,
    title=\textbf{Tier E --- Extension},
    fonttitle=\bfseries, arc=2mm, left=3mm, right=3mm, top=2mm, bottom=2mm]

\textbf{Comparing $f(x) = \log_2 x$ and $h(x) = 3\log_2 x$.}

\begin{enumerate}[leftmargin=1.8em, itemsep=10pt]

\item Build a table for both functions for $x = 1, 2, 4, 8, 16$:

\vspace{4pt}
\begin{center}
\renewcommand{\arraystretch}{1.8}
\begin{tabular}{|c|c|c|}
\hline
\rowcolor{sky}
$x$ & $f(x) = \log_2 x$ & $h(x) = 3\log_2 x$ \\
\hline
$1$  & \blank{2cm} & \blank{2cm} \\
\hline
$2$  & \blank{2cm} & \blank{2cm} \\
\hline
$4$  & \blank{2cm} & \blank{2cm} \\
\hline
$8$  & \blank{2cm} & \blank{2cm} \\
\hline
$16$ & \blank{2cm} & \blank{2cm} \\
\hline
\end{tabular}
\end{center}

\item Verify the additive growth property for $h(x) = 3\log_2 x$: as $x$ doubles, $h(x)$
      increases by \blank{2cm}. Is this additive growth? Explain.

\writeline

\item For what input(s) do $f(x)$ and $h(x)$ have the same output? Show your work algebraically.

\vspace{4pt}
\writelines{3}

\item Describe how the vertical dilation by 3 changes the output values of $h(x)$ compared
      to $f(x)$, and explain whether $h(x)$ is still considered a logarithmic function.

\writelines{2}

\end{enumerate}
\end{tcolorbox}

\end{document}
